
\documentclass[11pt]{article}
\usepackage[margin=1in]{geometry}
\usepackage{amsmath,amssymb,booktabs,hyperref,graphicx,parskip,enumitem}
\usepackage{xcolor}
\hypersetup{colorlinks=true,linkcolor=blue!50!black,urlcolor=blue!50!black}
\setlist{nosep}

\title{Replication Report: Molten Salt Reactor Neutronics and Fuel Cycle Modeling and Simulation with SCALE (MSRE Depletion)\\[2pt]\large OSTI 2217719 \textbar{} Coverage 8/10, Agreement 8/10}
\author{Rick Stevens \and Ollie (AI Assistant)\\\small Argonne National Laboratory}
\date{2026-04-27}

\begin{document}\maketitle

\section{Source paper}
\begin{itemize}
\item \textbf{Title:} Molten Salt Reactor Neutronics and Fuel Cycle Modeling and Simulation with SCALE (MSRE Depletion)
\item \textbf{Authors:} Betzler, Powers, Worrall
\item \textbf{Year:} 2017
\item \textbf{Domain:} Nuclear Engineering / Reactor Physics
\item \textbf{Identifier:} OSTI 2217719
\end{itemize}


\section{Replication objective}
The central claim of the paper concerns nuclear engineering / reactor physics. Our objective was to reproduce the headline numerical/qualitative results within the constraints of the AI-driven Replication Project (24--72\,h, commodity GPU/CPU compute, no author contact). A replication plan is archived at \texttt{2217719-SCALE-depletion-capabilities-for-molten-salt/replication\_plan.pdf}.

\section{Methodology}
We followed the standard Replication Project workflow: ingest the source PDF, extract method and data pointers, draft a replication plan, attempt the smallest end-to-end pipeline that produces a comparable number, then scale up where compute and data permit. Tools used in this replication are catalogued in the meta-paper's computational tools inventory (see \texttt{COMPUTATIONAL\_TOOLS\_INVENTORY.md}).



\subsection*{Paper-directory README excerpt}
\small
\par\medskip\textbf{SCALE depletion capabilities for molten salt reactors and other liquid-fueled systems}\par
\par$\bullet$\ **OSTI ID:** 2217719
\par$\bullet$\ **Rank:** \#5
\par$\bullet$\ **Replication Score:** 10/10

\par\medskip\textbf{\# Why This Paper}\par
This paper presents a computational methodology for simulating reactor depletion and material flows, with structured simulation outputs and clear equations, making it ideal for AI-driven replication and benchmarking.

\par\medskip\textbf{\# Replication Plan}\par
AI would implement the described ODEs and material flow models using SCALE input parameters, generate quantitative outputs for nuclide inventories and reactor metrics, and validate results against published test cases.

\par\medskip\textbf{\# Status}\par
\par$\bullet$\ [ ] Paper reviewed
\par$\bullet$\ [ ] Data identified
\par$\bullet$\ [ ] Code implemented
\par$\bullet$\ [ ] Results validated
\normalsize

The replication was driven by an OpenClaw agent loop (Claude Opus / GPT-5 class models) issuing shell, Python, and (where applicable) GPU jobs. Compute usage and wall-time for individual papers are summarised in the meta-paper's resource table; not separately recorded for this paper unless noted in the Tier-Lift artefact. Sub-agents were spawned for replication-plan generation and (for papers in batches 1--4) for tier-lift extension runs.

\section{What was reproduced}
\begin{itemize}
\item All previous Phase-2 items (MSRE 2D core, OpenMC depletion 375 d, k\_eff agreement)
\item Phase-1 three-mixture verification (paper Figs 3-6) via explicit Bateman ODEs
\item Pa-233 mass redistribution Mix1→Mix2,Mix3
\item U-233 buildup in waste mixtures from Pa decay
\item Nd-148 mass evolution under high-rate removal
\item Analytical I-135/Xe-135 equilibrium with online removal
\item Δρ\_Xe ≈ -1684 pcm with MSRE removal vs -3124 pcm without (Q1 followon)
\end{itemize}

\section{What was missing or incomplete}
\begin{itemize}
\item Full 3D SCALE/TRITON deterministic run
\item Full \textasciitilde{}900-day operational history with U-235 makeup
\item Broader fuel-cycle parameter sweeps
\item Spectral comparisons (fast/thermal flux distributions)
\item Temperature feedback / void coefficient studies
\end{itemize}
The dominant blocker categories for this paper, mapped onto the meta-paper's 9-category friction taxonomy (\S4), are listed in \S\ref{sec:friction} below.

\section{Quantitative agreement}
Per-quantity numerical comparisons are not separately tabulated in the structured evaluation record for this paper beyond the agreement rationale below; where the rationale references specific numbers, they are reproduced verbatim. A full numerical comparison table is recorded in the per-replication artefacts (see \S\ref{sec:repro}). For papers below 8/8, the agreement rationale identifies the specific quantities that diverge.

\paragraph{Agreement rationale.} All previous OpenMC Monte-Carlo depletion agreement preserved (k\_eff = 1.165 vs paper \textasciitilde{}1.16). Three-mixture Bateman test: Pa/U/Nd time evolution matches paper Figs 3-6 qualitatively (correct kinetics, correct asymptotic mass redistribution between mixtures). Xe-poisoning result is in line with literature MSRE values (\textasciitilde{}\$2 worth of reactivity worth of online removal benefit — operating reports of MSRE).

\section{Score and friction-point tagging}\label{sec:friction}
\begin{itemize}
\item \textbf{Coverage:} 8/10 --- Original c=6 lifted to 8 by (a) integrating the previously-completed Phase-1 three-mixture verification test (paper Figs 3-6: Pa-233, U-233, Nd-148 evolution across 3 mixtures with continuous removal streams, solved via explicit Bateman ODEs) into the replication dir, and (b) adding an analytical Xe-135 reactivity-sensitivity study (follow-on Q1) that solves the I-135/Xe-135 equilibrium pair as a function of online removal rate λ\_rem. Result: at MSRE λ\_rem = 4.067e-5/s, Δρ\_Xe = -1684 pcm; without removal Δρ\_Xe = -3124 pcm; benefit Δρ ≈ 1440 pcm (\textasciitilde{}\$2 reactivity, consistent with MSRE operating record). Local sensitivity dρ/d(ln λ\_rem) ≈ 776 pcm per e-fold. Still missing: full \textasciitilde{}900-day operational history with U-235 makeup, broader fuel-cycle parameter sweeps.
\item \textbf{Agreement:} 8/10 --- All previous OpenMC Monte-Carlo depletion agreement preserved (k\_eff = 1.165 vs paper \textasciitilde{}1.16). Three-mixture Bateman test: Pa/U/Nd time evolution matches paper Figs 3-6 qualitatively (correct kinetics, correct asymptotic mass redistribution between mixtures). Xe-poisoning result is in line with literature MSRE values (\textasciitilde{}\$2 worth of reactivity worth of online removal benefit — operating reports of MSRE).
\end{itemize}

\noindent\textbf{Friction points encountered (taxonomy tags):}
\begin{itemize}
\item None recorded.
\end{itemize}

\section{Follow-on questions}
\begin{itemize}
\item Q2/Q3/Q4/Q5 from prior entry remain open
\end{itemize}

\section{Reproducibility pointers}\label{sec:repro}
\begin{itemize}
\item \textbf{Paper directory:} \texttt{2217719-SCALE-depletion-capabilities-for-molten-salt}
\item \textbf{Replication artefacts:}
\begin{itemize}
\item Replication artifacts in paper directory.
\end{itemize}
\item \textbf{Replication-dir hint from JSONL:} \texttt{\textasciitilde{}/Dropbox/REPLICATE-PROJECT/2217719-scale-msr/replication/}
\item \textbf{Status:} tier\_lift\_2026\_04\_26
\end{itemize}

To re-run, change directory into the paper directory above and consult its \texttt{README.md} for the entry-point script. Most replications follow the convention \texttt{replication/run\_*.sh} or \texttt{replication/<subdir>/run.py}.

\end{document}
